% --- Tabla: Relación entre error, defecto y fallo (Actividad 1) ---
\begin{table}[H]
  \centering
  \small
  \caption{Identificación de error, defecto y fallo en casos de prueba y en el proyecto integrador.}
  \label{tab:pruebas_error_defecto}
  \begin{tabularx}{\linewidth}{|L{3.5cm}|Y|Y|Y|}
    \hline
    \textbf{Situación} & \textbf{Error (programador)} & \textbf{Defecto (en el software)} & \textbf{Fallo (usuario)} \\ \hline
    1. En el módulo de transferencias, el sistema rechaza la operación cuando el monto es exactamente igual al saldo.
      & El programador confundió el operador $>$ con $\geq$.
      & Condición \texttt{saldo > monto} en la función \texttt{transferir()}.
      & Mensaje «saldo insuficiente» aunque el saldo alcanza. \\ \hline
    2. Un formulario de registro acepta el correo «juan@» y luego no se envía la confirmación.
      & El programador no validó el formato del correo (omitió comprobar la estructura \texttt{usuario@dominio}).
      & El campo correo del formulario/controlador de registro no tiene validación de formato.
      & El registro se guarda con «juan@», pero nunca llega el correo de confirmación. \\ \hline
    3. Un reporte de ventas muestra el total del mes anterior al seleccionar el mes actual.
      & El programador usó el mes anterior (mes $-$ 1) o ignoró el parámetro de mes al armar la consulta.
      & El filtro de fechas de la consulta del reporte toma un rango de mes incorrecto.
      & Al elegir el mes actual, el reporte muestra el total del mes anterior. \\ \hline
    4. \textbf{Caso propio del proyecto:} en el POS se agrega dos veces el mismo juguete (2 y 3 unidades) cuando solo quedan 4 en stock, y la venta se acepta.
      & El programador validó la cantidad de cada línea por separado en lugar de la cantidad total pedida de cada juguete.
      & En \texttt{VentaController::\allowbreak store()} no se agrupaban los ítems por \texttt{id\_producto} antes de comparar con el stock (RF-03).
      & El ticket se genera por 5 unidades y, al cobrar en caja, el inventario queda en negativo o el cobro se rechaza con el cliente presente. \\ \hline
  \end{tabularx}
  \par\vspace{0.2cm}
  {\raggedright \small \emph{Nota.} Elaboración propia (2026). Análisis de causalidad entre error humano, defecto en código y fallo perceptible, incluyendo un caso real detectado en ``Mundo del Juguete''.\par}
\end{table}
